\documentclass[a4paper,12pt]{article}
\usepackage{tikz}
\usepackage{graphicx}
\usepackage{soul}
\usepackage{multirow}
\usepackage{float}
\usepackage{array,booktabs,arydshln,xcolor}
\usepackage{amsthm}
\usepackage{amsmath,amssymb}
\usepackage{pifont}
\newcommand{\xmark}{\ding{53}}
\usepackage[framed,numbered,autolinebreaks,useliterate]{mcode}
\usepackage{caption}
\usepackage{subcaption}
\usepackage[margin=2.5cm]{geometry}
\usepackage{adjustbox}
\usepackage{diagbox}
\usepackage{systeme} 
\usepackage[parfill]{parskip}
\usepackage[shortlabels]{enumitem}
\usepackage{listings}
\usepackage{fancyhdr}
\usepackage{bbm}
\usepackage[backend=bibtex, style=authoryear]{biblatex}
%\usepackage[style=authoryear]{biblatex}
%\usepackage[authoryear]{natbib}
\usepackage{tabularx,ragged2e}
\newcolumntype{C}[1]{>{\Centering}m{#1}}
\renewcommand\tabularxcolumn[1]{C{#1}}
\usepackage{makecell}
\usepackage[T1]{fontenc} 
\usepackage{fouriernc}
\usepackage{mathtools}
\usepackage{ marvosym }
\usepackage{multicol}
\usepackage[toc,page]{appendix}
\newcommand\blitz{\textup{\Lightning}}
\usepackage[graphicx]{realboxes}
\usepackage{titlesec}
\setcounter{secnumdepth}{4}
\titleformat{\paragraph}
{\normalfont\normalsize\bfseries}{\theparagraph}{1em}{}
\titlespacing*{\paragraph}
{0pt}{3.25ex plus 1ex minus .2ex}{1.5ex plus .2ex}
\usepackage{hyperref}
\newlist{todolist}{itemize}{2}
\setlist[todolist]{label=$\square$}
\newcommand{\cmark}{\ding{51}}%
\newcommand{\done}{\rlap{$\square$}{\raisebox{2pt}{\large\hspace{1pt}\cmark}}%
\hspace{-2.5pt}}
\newcommand{\wontfix}{\rlap{$\square$}{\large\hspace{1pt}\xmark}}
\theoremstyle{plain}
\newtheorem{theorem}{Theorem}[section]
\newtheorem{lemma}[theorem]{Lemma}
\newtheorem{corollary}[theorem]{Corollary}
\newtheorem{prop}[theorem]{Proposition}
\newtheorem{conj}[theorem]{Conjecture}

\theoremstyle{definition}
\newtheorem{definition}[theorem]{Definition}
\newtheorem{example}[theorem]{Example}
\newtheorem{examples}[theorem]{Examples}

\newtheorem{question}[theorem]{Question}
\newtheorem{questions}[theorem]{Questions}

%\theoremstyle{remark}
\newtheorem{remark}[theorem]{Remark}

\newtheoremstyle{TheoremNum}
        {\topsep}{\topsep}              %%% space between body and thm
        {\itshape}                      %%% Thm body font
        {}                              %%% Indent amount (empty = no indent)
        {\bfseries}                     %%% Thm head font
        {.}                             %%% Punctuation after thm head
        { }                             %%% Space after thm head
        {\thmname{#1}\thmnote{ \bfseries #3}}%%% Thm head spec
    \theoremstyle{TheoremNum}
    \newtheorem{theoremn}{Theorem}


\newtheoremstyle{CorollaryNum}
        {\topsep}{\topsep}              %%% space between body and thm
        {\itshape}                      %%% Thm body font
        {}                              %%% Indent amount (empty = no indent)
        {\bfseries}                     %%% Thm head font
        {.}                             %%% Punctuation after thm head
        { }                             %%% Space after thm head
        {\thmname{#1}\thmnote{ \bfseries #3}}%%% Thm head spec
    \theoremstyle{CorollaryNum}
    \newtheorem{corollaryn}{Corollary}

\newtheoremstyle{PropositionNum}
        {\topsep}{\topsep}              %%% space between body and thm
        {\itshape}                      %%% Thm body font
        {}                              %%% Indent amount (empty = no indent)
        {\bfseries}                     %%% Thm head font
        {.}                             %%% Punctuation after thm head
        { }                             %%% Space after thm head
        {\thmname{#1}\thmnote{ \bfseries #3}}%%% Thm head spec
    \theoremstyle{PropositionNum}
    \newtheorem{propn}{Proposition}

\newtheorem{conjecture}[theorem]{Conjecture}

    \newtheoremstyle{ConjectureNum}
        {\topsep}{\topsep}              %%% space between body and thm
        {\itshape}                      %%% Thm body font
        {}                              %%% Indent amount (empty = no indent)
        {\bfseries}                     %%% Thm head font
        {.}                             %%% Punctuation after thm head
        { }                             %%% Space after thm head
        {\thmname{#1}\thmnote{ \bfseries #3}}%%% Thm head spec
    \theoremstyle{TheoremNum}
    \newtheorem{conjn}{Conjecture}
    
\newcommand{\argmin}{\arg\!\min}
\newcommand{\argmax}{\arg\!\max}
    
\addbibresource{main.bib}
\tolerance=1
\emergencystretch=\maxdimen
\hyphenpenalty=10000
\hbadness=10000
\linespread{1}

%opening
\title{Divorce-induced health decline and the impact on mortality risk\thanks{The authors are grateful to the DFG (German Research Foundation) for financial support (Project number: 544381616).}}

\author{Kar Man Tan\thanks{Corresponding author: \href{mailto:karman.tan@bib.bund.de}{karman.tan@bib.bund.de}. Federal Institute for Population Research (BiB), Friedrich-Ebert-Allee 4, 65185 Wiesbaden, Germany}
\and
Pavel Grigoriev\thanks{Federal Institute for Population Research (BiB)}}
% Pavel Grigoriev\thanks{\href{mailto:pavel.grigoriev@bib.bund.de}{pavel.grigoriev@bib.bund.de}. Federal Institute for Population Research (BiB)}}

\begin{document}

\maketitle

\begin{abstract}
Divorce is linked to elevated mortality. However, the causal mechanisms underlying this relationship are complex and remain insufficiently identified. Health is one such mechanism and is the focus of this paper. Using linked administrative data from the German Pension Insurance Fund (DRV), we examine whether divorce increases the risk of an initial musculoskeletal rehabilitation, and how much of the divorce-mortality gap runs through this channel. Rehabilitation is administered by the DRV and is granted only once a musculoskeletal condition has become chronic and threatens the ability to work. Our outcome measures the crossing of an institutional threshold rather than the onset of disease. We build annual risk sets of individuals experiencing a first (observable) divorce and individuals who have not yet divorced in that year and then use propensity score matching to generate comparable matches. \textcolor{red}{Results: Divorce effect on cumulative incidence of musculoskeletal rehab, divorce effect on mortality, results from decomposition}
% \textit{(JEL D41, D42, D43, D81, D82, D83, D84, D85, D86)}
\end{abstract}

\textbf{Keywords:} divorce, health, rehabilitation, mortality

\textbf{JEL Classification:} H75, I12, J12

%[SUBMITTED TO \textbf{JOURNAL OF RISK AND INSURANCE} AND IN REVIEW]

\newpage

\section{Introduction}

The link between divorce and mortality has been widely discussed in social epidemiological studies. Register-based studies have used Finnish data (\citeauthor{metsa2013short}, \citeyear{metsa2013short}), Norwegian data (\citeauthor{berntsen2012relationship}, \citeyear{berntsen2012relationship}) and  English/Welsh data (\citeauthor{grundy2010marital}, \citeyear{grundy2010marital}). Until recently, Germany has been a difficult setting for this question. \citeauthor{brockmann2004love} (\citeyear{brockmann2004love}) is the only study to the best of our knowledge which explicitly analyzes the link between divorce and mortality although the dataset used in the study is the SOEP\footnote{German Socioeconomic Panel} which is not an administrative dataset. Since 2011, it is possible to link different administrative datasets from the German Pension Insurance Fund (DRV) using a unique identifier, allowing us to access full income biographies, divorce information\footnote{This information is available due to the law of splitting pension points.}, rehabilitation history and death. The availability of this dataset led to the conception of the project \textit{Divorce and Diagnosis (DiDi): Register-based Analysis on the Role of Divorce in Health and Employment Dynamics}\footnote{This project is a collaboration between the Federal Institute for Population Research, Hertie School, Charit\'{e} - Universit\"{a}tsmedizin Berlin and the DRV.}, which constitutes our study. 

Our study focuses on the interaction between divorce, health and mortality. While there can be a direct link between divorce and mortality if divorce results in suicide (\citeauthor{kposowa2000marital}, \citeyear{kposowa2000marital}; \citeauthor{stevenson2006bargaining}, \citeyear{stevenson2006bargaining}) or homicide, causal evidence remains scarce when the link is indirect. Hence, we chose health as it is one possible indirect link tying together divorce and mortality. There can, however, still be threats to identification as health is not only a consequence but also a cause of divorce  (\citeauthor*{torvik2015health}, \citeyear{torvik2015health}). Furthermore, the individuals who choose to divorce may not necessarily be random. Many of the register-based studies make use of hazard models to estimate excess mortality risk after divorce and controls for sociodemographic and socioeconomic variables. For Germany, \citeauthor{brockmann2004love} (\citeyear{brockmann2004love}) account for health by using the self-reported health measure from the German SOEP and control for variables associated with marital status to find that divorced mean have 51\% higher mortality than married men while there is no significant effect for women. \citeauthor{leopold2018gender} (\citeyear{leopold2018gender}) presents a very careful multi-outcome treatment of the gendered consequences of divorce over time, but it is not designed to control for the health that occurs before from the health consequences which follow the divorce. In our paper, we control for pre-existing health trends and other factors which occur before the divorce by using propensity score matching to answer the following questions

\begin{itemize}
    \item[\textbf{RQ1:}] Does divorce raise the risk of musculoskeletal rehabilitation\footnote{Musculoskeletal diseases are the most common cuase of lost work time in Germany.} among divorced and non-divorced individuals?
    \item[\textbf{RQ2:}] Is some of the excess mortality caused by divorce occurring because divorce caused additional musculoskeletal deterioration?
\end{itemize}

For context, the DRV only grants rehabilitation once the musculoskeletal condition is chronic to the point that it threatens working ability (\citeauthor{drv2005leitlinien}, \citeyear{drv2005leitlinien}). We do not argue that divorce causes musculoskeletal disease but rather whether divorce moves individuals who carry a latent vulnerability across the legal threshold where rehabilitation becomes necessary. Please see Section 2.3 for a more detailed discussion. 

\textcolor{red}{Results... maybe policy implications?}

Section 2 provides a literature review, Section 3 introduces the data and methodology, Section 4 presents the results, Section 5 discusses the implications and Section 6 concludes.

\section{Literature review}
\subsection{Link between divorce and mortality}
The link between divorce and mortality risk is widely documented in sociology and epidemiology. In a meta-analysis, \citeauthor*{shor2012meta} (\citeyear{shor2012meta}) run meta-regressions on 625 mortality risk estimates over 104 studies and find that the mean hazard ratio was 1.30 for divorced individuals. The hazard ratio is slightly higher for men at 1.37 and lower for women at 1.22, although the difference between genders decrease with age. However, heterogeneity on the effect on death exists depending on the time since divorce. \citeauthor{metsa2013divorce} (\citeyear{metsa2013divorce}) finds that excess mortality among Finnish men, tied to accidental, violent or alcohol-related deaths, is highest immediately after divorce and declines over time. On the other hand, \citeauthor{berntsen2012relationship} (\citeyear{berntsen2012relationship}) finds that excess mortality among Norwegian men increases with time since divorce and remarriage can change the magnitude of the effect. In Germany, the evidence is more limited. \citeauthor{brockmann2004love} (\citeyear{brockmann2004love}) use event-history models which account for marital transitions and duration in each marital status. They find evidence of selection into marriage and that mortality risk is lower when the marriage duration is longer. 



Most of the previous studies establish a robust association, however, causal evidence remains scarce as divorce is not random. \citeauthor{sbarra2009divorce} (\citeyear{sbarra2009divorce}) follow 1376 individuals in the Charleston Heart Study for forty years and show that the definition of the exposure determines the size of the effect. Individuals who are continuously separated or divorced have a higher hazard ratio (1.66) compared to individuals who have ever been separated or divorced (0.98). \citeauthor{wiborg2025partnership} (\citeyear{wiborg2025partnership}) initially find raised hazard ratios for divorced men (1.59) and women (1.47) using Norwegian data. After using twin-pairs to account for genetic and early developmental and environmental factors, the relationship between divorce and mortality disappears.

\subsection{Divorce and health}

The effect of divorce on health is difficult to isolate due to the possibility of reverse causality. \citeauthor*{joung1998longitudinal} (\citeyear{joung1998longitudinal}), using Dutch data, find that married individuals who had more subject health complaints were 1.5 times more likely to divorce, while individuals who reported multiple chronic conditions were twice as likely to divorce. \citeauthor{torvik2015health} (\citeyear{torvik2015health}) corroborates \citeauthor*{joung1998longitudinal} (\citeyear{joung1998longitudinal})'s findings with Norwegian data. Poor health behaviors including smoking, drinking, and lack of exercise in addition to mental distress predict subsequent divorce. \citeauthor{metsa2013divorce} (\citeyear{metsa2013divorce}), with Finnish data, find that psychotropic medication use among individuals who divorce is already elevated five years before hand and peaks at around six to nine months just before divorce. \citeauthor{wade2004marital} (\citeyear{wade2004marital}) find worse mental health among those divorcing or divorced than among the married using British data. Generally, poor health tends to precede divorce.

To deal with health selection, \citeauthor{kohn2014effect} (\citeyear{kohn2014effect}) estimates a dynamic model that accounts for the persistence of health and how it interacts with relationship status using 18 waves of the British Household Panel Survey. The result is that divorce has a marginally significant effect on the health of men younger than 45. \citeauthor*{sbarra2014marital} (\citeyear{sbarra2014marital}) also finds only a marginal effect of divorce on health. After matching individuals on divorce probability, \citeauthor*{sbarra2014marital} (\citeyear{sbarra2014marital}) finds that only individuals who were already depressed before divorce experienced major depressive episodes following divorce. With respect to mental health, \citeauthor{bjorkenstam2013risk} (\citeyear{bjorkenstam2013risk}) find that divorced individuals have thrice the risk of psychiatric inpatient care compared with married individuals in Sweden. \citeauthor{meulman2023estimating} (\citeyear{meulman2023estimating}) find similar results in the Netherlands where divorced individuals had higher health expenditures dedicated to mental health care. In terms of physical health, the evidence is heterogeneous. \citeauthor{zulkarnain2019divorce} (\citeyear{zulkarnain2019divorce}) find that women who divorce at older ages tend to have more physical changes than men. \citeauthor{leopold2018gender} (\citeyear{leopold2018gender}), on the other hand, find that men are more vulnerable in the short run in terms of well-being although the gender differences narrow in the medium run. \citeauthor{pellon2024divorce} (\citeyear{pellon2024divorce}), in a meta-analysis across 94 studies find that divorced individuals report worse self-rated physical conditions with higher complaints of musculoskeletal, cardiovascular and cerebrovascular conditions although the effect is not significant.


\subsection{Divorce and musculoskeletal disease}

Evidence on the pathway from divorce to musculoskeletal impairment is rather limited and ambiguous. However, we chose to focus on musculoskeletal diseases because they are a major source of disability and work interruption in Germany\footnote{Please see the section on \nameref{app:rehab_germany} in the Appendix.}. The DRV Reha-Atlas (\citeauthor{drv2024rehaatlas}, \citeyear{drv2024rehaatlas}) reports of all completed medical rehabilitations, orthopedic conditions are the most common, accounting for 35.6\% of rehabilitations among women and 35.5\% among men. \citeauthor{baua2026arbeitswelt} (\citeyear{baua2026arbeitswelt}) reports that musculoskeletal diseases account for 21\% of certified incapacity days in 2024, which is the largest diagnostic group. 

\citeauthor{sjogren2009epidemiology} (\citeyear{sjogren2009epidemiology}), using Danish data, find that chronic pain is more common among individuals who are divorced, separated or widowed, where musculoskeletal diseases accounted for 66.8\% of the cause of chronic pain. However, \citeauthor{choi2026marital} (\citeyear{choi2026marital}), using the Health and Retirement Survey find that individuals who transition from marriage to divorce tend to have a lower probability of high-impact pain. On the other hand, \citeauthor{cimas2018chronic} (\citeyear{choi2026marital}) using SHARE\footnote{Survey of Health, Aging and Retirement in Europe} data, find that there are gender differences. While divorced women have a higher incidence of chronic musculoskeletal pain, men have a lower prevalence. Meanwhile, \citeauthor{andorsen2017prospective} (\citeyear{andorsen2017prospective}) finds that there is no significant effect between divorce and musculoskeletal complaints using Norwegian data.

While the impact of divorce on musculoskeletal diseases is ambiguous, the evidence on the relationship between social support and the consequences of musculoskeletal diseases is more homogeneous. \citeauthor*{prang2015recovery} (\citeyear{prang2015recovery}) finds that individuals who have support from family friends and neighbors tend to have less persistent pain associated with musculoskeletal diseases, while widowed, separated or divorced individuals, in the absence of a network, tend to have a higher probability of persistent pain. \citeauthor{woods2005work} (\citeyear{woods2005work}) corroborates this and finds that a lack of support is associated with higher musculoskeletal morbidity, absence from work and subsequently, failure to return to work. \citeauthor{reese2010pain} (\citeyear{reese2010pain}) finds that the quality of marriages matter, i.e. individuals with arthritis in problematic marriages tend to have more pain complaints and disability compared to individuals in well-adjusted marriages. This could signify that a decline in the relationship quality prior to divorce leads to poorer support, hence more complaints of pain. \citeauthor*{couch2015long} (\citeyear{couch2015long}) and \citeauthor*{tamborini2016work} (\citeyear{tamborini2016work}) find that divorce is associated with subsequent health threatening the ability to work, especially among individuals who do not remarry, although they do not differentiate by types of diseases. 


\subsection{Musculoskeletal disease and death}

Musculoskeletal conditions are rarely direct causes of death in the way that cardiovascular diseases are. They often operate indirectly, for example through inflammation affecting the cardiovascular disease risk. Rheumatoid arthritis and chronic inflammatory activity can contribute to oxidative stress, dysfunction of blood vessels, arterial build-up and other cardiovascular risks (\citeauthor{england2018increased}, \citeyear{england2018increased}). Cardiovascular mortality among people with rheumatoid arthritis is about 60\% higher than the general population (\citeauthor{meulman2023estimating}, \citeyear{meune2009trends}). For non-inflammatory conditions, musculoskeletal conditions causes pain and restricted movement, which then makes people avoid physical activity. This reduces cardiovascular health and fitness, which then potentially increases the risk of comorbidities. \citeauthor{nuesch2011all} (\citeyear{nuesch2011all}) finds that English patients with knee or  hip osteoarthritis tend to have higher all-cause and cardiovascular mortality. \citeauthor{hawker2014allcause} (\citeyear{hawker2014allcause}) finds a similar result with Canadian patients. 

\citeauthor{jansson2012sickness} (\citeyear{jansson2012sickness}) is more related to our study. Using Swedish register data, the authors compared individuals who were absent from work due to musculoskeletal diseases with people who were absent from work due to other diseases and people who were not absent from work. They find that among three categories of musculoskeletal diseases, the all-cause morbidity was 1.3 to 1.5 times higher. Given that musculoskeletal diseases account for the highest share of rehabilitations facilitated by the DRV, it is important that rehabilitation can encourage return to work and by extension, reduce the all-cause mortality attributed to musculoskeletal diseases. \citeauthor{bethge2019effects} (\citeyear{bethge2019effects}) finds that work-related medical rehabilitation tends to increase the return to work rate by 5.8 percentage points and shortens the median time to return by 9.46 days compared to conventional medical rehabilitation. In our study, the first completed DRV-facilitated musculoskeletal rehabilitation should be interpreted as an indicator of serious, work-relevant impairment and also receipt of potentially beneficial treatment. 


%Musculoskeletal conditions are rarely direct causes of death in the way that ischaemic heart disease or cancer are. They are slow-burning, and raise mortality indirectly through disability, comorbidity accumulated under a more sedentary life, pain, falls and fracture (Williams, et al., 2018). Chronic musculoskeletal pain is in turn linked to smoking, depression and anxiety, each of which carries its own mortality risk (Chen, et al., 2021). Sickness absence itself predicts mortality (Vahtera, Pentti and Kivimaeki, 2004; Bambra and Norman, 2006; Mittendorfer-Rutz, et al., 2012), which is why rehabilitation entry is a defensible mediator to test - but also why its sign is not obvious ex ante, since entering rehabilitation is simultaneously a marker of impairment and a treatment for it.

\section{Data and methodology}
\subsection{Data}
The Research Data Center of the DRV allows linkages of different datasets starting from the reporting year 2011 through a unique identifier (Forschungsdatenzentrum Rentenversicherung, 2025). The datasets we use are AKVS\footnote{Active Insured Persons}, LTVS\footnote{Latent Insured Persons}, RTBN\footnote{Insured Persons’ Pension Portfolio}, RTWF\footnote{Termination of Pensions due to Death}, RSD\footnote{Completed Rehabilitation in the Insurance History} and VA\footnote{Pension Equalization Statistics}. The time period covered by the data is 2011 to 2023.

The AKVS, LTVS and RTBN are merged by the DRV to form the BASE dataset. The AKVS dataset contains individuals who are still actively contributing to the pension fund hence it generally represents individuals who are employed. The LTVS dataset contains individuals who have paused pension fund contributions but have not yet retired (nor drawn from their pension fund). It may include individuals who are unemployed, undergoing further education, on parental leave, etc. The RTBN dataset contains individuals who are drawing from their pension fund, i.e. retired individuals. The merged dataset, BASE, allows us to map out the full life-course coverage of individuals. If we only considered only RTBN, the sample is biased towards individuals who have contributed and lived long enough to draw a pension, ignoring those individuals who have undergone rehabilitation or died before drawing a pension. Considering only AKVS ignores rehabilitation efforts and gray divorces. To round up the dataset, it is also important to consider LTVS which details what happens between employment spells and retirement.

We then supplement the merged dataset with the RTWF, RSD and the VA. The RTWF details pension terminations due to deaths. If the death information from the merged dataset is inconsistent with the information in the RTWF, the individuals (cases) are dropped \textcolor{red}{\{about 127 individuals are dropped\}}. The RSD dataset details the completed rehabilitations and include information such as the broad diagnosis group, the type of rehabilitation benefits applied for, start and end date of rehabilitation. The RSD dataset also includes individuals with withdrawn, transferred or rejected applications. We only include individuals whose applications were successful and who completed the rehabilitation spell. Furthermore, about 5\% of individuals undergo three or more rehabilitation spells per year. We also drop these individuals. We then supplement the merged dataset by including the additional rehabilitation information for the remaining individuals in the RSD dataset. The VA dataset includes information related to divorce and only contains divorced individuals. From the VA, we can obtain the date of divorce, pension supplements and deductions which result from divorce.


\subsection{Matching}

Since divorce is not randomly assigned, divorced and non-divorced individuals may differ systematically before divorce occurs. We therefore use 1:1 matching without replacement to construct a comparison group of individuals who are similar with respect to observed pre-treatment characteristics. The causal interpretation relies on conditional exchangeability (\citeauthor{rosenbaum1983central}, \citeyear{rosenbaum1983central}). This means that, conditional on the observed characteristics used in matching, individuals who divorce and comparable individuals who do not divorce are assumed to have had similar potential health trajectories in the absence of divorce. In other words, it is assumed that there are no remaining unobserved factors that jointly affect the probability of divorce and the subsequent health outcomes. However, the administrative data may not fully capture factors such as relationship conflict or social support.

We define treatment as the first observed divorce from the earliest observed marriage spell. We limit the treatment cohorts to divorces observed between 2012 and 2018 to ensure that each cohort can be observed for at least five years after divorce. The year of divorce is fixed as $t_0$. For each treatment year, we form a separate within-year risk set. For treated individuals whose first observed divorce occurs in $t_0$, eligible controls are individuals who have not yet experienced a divorce by that year. Controls may divorce later. We do not restrict the control group to individuals who never divorce because this would define control eligibility using information that occurs after $t_0$. This means that it is possible that an individual who was a control before their own divorce can be counted as treated afterwards. Individuals must be alive at $t_0$ and must have an observed record in $t_{0}-1$. We exclude individuals with a musculoskeletal rehabilitation up to and including $t_{0}$. They may previously have received rehabilitation for other conditions. The resulting sample therefore consists of individuals who are at risk of their first observed musculoskeletal rehabilitation after $t_0$.
 
We require treated and control individuals to be matched exactly on \texttt{gender}, five-year \texttt{age\_band} and whether any rehabilitation has occurred by $t_{0}-1$, i.e. \texttt{ever\_had\_any\_rehab\_to\_date}. Then, we estimate a logistic propensity score model separately for each treatment year. The propensity score represents the estimated probability of divorce in $t_0$, conditional on observed characteristics measured before divorce at $t_{0}-1$. The numeric characteristics we include are \texttt{age}, \texttt{income}, \texttt{pension\_income}, \texttt{cumulative\_rehabs\_to\_date}, \texttt{cumulative\_mental\_rehabs\_to\_date}, \texttt{years\_since\_last\_rehab}, \texttt{full\_time\_contribution\_periods}\footnote{Periods during which pension contributions were recorded for full-time work.} and \texttt{part\_time\_contribution\_periods}\footnote{Periods during which pension contributions were recorded for part-time work}. The binary characteristics we include are \texttt{ever\_had\_mental\_rehab\_to\_date}, \texttt{unsuccessful\_rehab\_application\_this\_year}, \texttt{missing\_full\_time\_contribution} and \texttt{missing\_part\_time\_contribution}. The categorical variables included are \texttt{federal\_state}, \texttt{broad\_occupational\_category}\footnote{Classification of Occupations 1988 created by the German Federal Employment Agency. We used the first digit which signifies occupational area.}, \texttt{income\_source}\footnote{This variable has the following categories regarding the source of income: 1) income only from working, income only from pensions, income from earnings and pensions, no income from earnings or pensions. The latter two categories are uncommon.} and \texttt{marital\_status}\footnote{Annual marital status is reconstructed from marriage spell records. In a few cases where an individual is married for exactly half the year, status is assigned by an unseeded random draw at the data processing stage. This can affect the variable used for propensity score balancing.}.

We treat missingness in earnings, pension income and contribution periods as potentially informative. In particular, observed earnings with missing pension income generally indicate active employment, whereas observed pension income with missing earnings generally indicates pension receipt. We retain the observed earnings and pension values and include a categorical measure describing the joint observability of the two income sources. We also include separate indicators for missing contribution-period information. This allows the matching model to use both the recorded values and the information contained in their missingness patterns.

Within each treatment year and exact-match group, we use optimal matching. The matching procedure first maximizes the number of treated individuals who can be matched within the caliper and then chooses the combination of pairs that minimizes the total propensity score distance. We limit matches to a maximum absolute propensity score difference of 0.01. Matching is conducted without replacement, meaning that each treated individual and each control can appear in only one matched pair. To accommodate the large annual control pools, the propensity score model is fitted using all treated individuals and, where necessary, a reproducibly selected sample of controls, with the total estimation sample limited to 500,000 observations\footnote{Please see the section on \nameref{app:control_limitation_500_000} in the Appendix}. The resulting propensity scores are corrected for the control-sampling fraction, and the fitted model is then used to calculate propensity scores for every eligible treated individual and control. The computational limit therefore affects how the propensity model is estimated but does not remove controls from the matching pool.

\textcolor{red}{We assess covariate balance before and after matching using standardized mean differences. Balance is assessed separately for continuous variables, binary variables, each level of categorical variables and indicators of missingness. We present these results in Figure blablabla. An absolute standardized mean difference below 0.10 is treated as evidence of adequate balance on the corresponding observed characteristic. Where residual imbalance above this threshold remains after matching, the affected variables are additionally included in the outcome model. This adjustment is intended to address limited remaining imbalance and is not used as a substitute for unsuccessful matching.}


\textcolor{red}{Using the matched sample, we estimate the average treatment effect among the matched treated individuals. We examine first entry into musculoskeletal rehabilitation, time to first musculoskeletal rehabilitation, rehabilitation duration and all-cause mortality after divorce. Musculoskeletal rehabilitation is examined over a five-year follow-up period, with death treated as a competing event. Mortality is followed for up to ten years where permitted by the end of the observation period in 2023. The amount of available mortality follow-up therefore varies across treatment cohorts.}


\subsection{Identification}

For identification, we focus on the average treatment effect among the matched treated individuals. 
\subsubsection{RQ1: Musculoskeletal rehabilitation}

We follow each matched individual from $t_0+1$ to $t_0+5$ and estimate a discrete-time hazard model:

\begin{align}
\begin{split}
 \lambda_{i,\tau}^{\text{rehab}} & = \text{Pr}(R_{i \tau} = 1 | \text{no musculoskeletal rehab before $\tau$, alive at $\tau$}) \\
 & = \Lambda_{\text{rehab}}\left(\alpha_\tau + \beta D_i + \mathbf{x}_i'\mathbf{\gamma}\right)
\end{split}
\tag{1}\label{eq:rq1_equation}
\end{align}

where $R_{i,\tau}$ is a dummy indicating individual $i$'s first entry into musculoskeletal rehabilitation in time in follow-up year $t_0 + \tau$ where $\tau \in \{1,...,5\}$, $\alpha_\tau$ are fixed effects, $D_i$ is a dummy indicating divorce at $t_0$ and $\mathbf{x}_i$ are other variables measured at $t_0 -1$ that we add depending on the specification. We estimate two specifications for Equation \ref{eq:rq1_equation}. In the first specification, we omit $\mathbf{x}_i$ and rely only on matching to deal with selection bias. In the second specification, we include variables whose absolute standardized mean difference after matching exceed 0.10 into $\mathbf{x}_i$. Standard errors are clustered at the level of the matched pair. 

The hazard odds ratio from Equation \ref{eq:rq1_equation} is difficult to interpret on its own as it does not show how frequently rehabilitation occurs, hence we also estimate the cumulative incidence of a first musculoskeletal rehabilitation over the 5-year follow up period using a discrete Aalen-Johansen estimator\footnote{This estimator recognizes that individuals who die can no longer enter rehabilitation. For example, suppose there are 100 people where 10 enter rehabilitation in $t_0 + 1$, 20 die before rehabilitation and 70 are still alive and have not entered rehabilitation. Suppose in $t_0 + 2$, 7 people out of the remaining 70 enter rehabilitation. If we treat death as censoring rather than as a competing event, the proportion surviving up to $t_0+1$ is $1-\frac{10}{100} = 0.9$, and up to $t_0+2$ is $1-\frac{7}{70} = 0.9$, $\therefore$ the cumulative rehabilitation probability is $1-(0.9\times 0.9) = 0.19$, which is an overestimation. With the Aalen-Johansen estimator, we consider the proportion which actually entered rehabilitation by each follow-up year before dying. In $t_0+1$, the proportion entering rehabilitation is 0.1, however, after $t_0+1$, the proportion remaining alive without having entered rehabilitation is $1-\frac{10+20}{100} = 0.7$. In $t_0+2$, $\frac{7}{70}=0.1$ is the share of people who enter rehabilitation, but since only 70\% of the original population was alive without having entered rehabilitation at the beginning of $t_0+2$, the contribution to the incidence is $0.7\times0.1=0.07$, hence the cumulative rehabilitation probability is $0.17$} (\citeauthor{aalen1978empirical}, \citeyear{aalen1978empirical}). We report four values at the the 5-year horizon, 1) cumulative rehabilitation risk among treated individuals, 2) cumulative rehabilitation risk among controls, 3) risk difference per 1000 individuals and 4) the corresponding risk ratio. Confidence intervals are obtained by resampling whole matched pairs. 

The biggest threat to the conditional exchangeability assumption is unobserved pre-divorce trends, which available variables may not fully capture, especially by only observing the year prior to the divorce. To address this, we first conducted an event study from $t_0 -3$ to $t_0+5$ by taking the difference in values of unsuccessful rehabilitation applications, earnings, pension income, missingness of pension income and prior rehabilitation activity between the treated and control in each matched pair and conducting a Wald test (\citeauthor{wald1943tests}, \citeyear{wald1943tests}) to determine whether the differences across time are significant for each variable. Following divergences in trends prior to $t_0$, we repeated the matching exercise with three lags, i.e. $t_0 -3, t_0-2, t_0-1$ instead of only $t_0 -1$. Since we require that individuals have an observed record, the lag-3 specification considers only the time period between 2014 and 2018. We also repeat the lag-1 specification for only 2014 to 2018.


\subsubsection{RQ2: Mortality}

We follow each individual from $t_0+1$ until $t_0+10$ or 2023, the last year of the sample, whichever occurs first and estimate a discrete-time hazard model similar to Equation \ref{eq:rq1_equation}:

\begin{align}
\begin{split}
    \lambda_{i,\tau}^{\text{death}} & = \text{Pr}(Y_{i,\tau}=1 | \text{alive at beginning of follow-up year } \tau)\\
    & = \Lambda_{\text{death}}\left(\alpha_\tau + \beta D_i + \mathbf{x}_i'\mathbf{\gamma}\right)
\end{split}
\tag{2}\label{eq:rq2_equation_basic}
\end{align}

where $Y_{i,\tau}$ is a dummy indicating that individual $i$ dies in follow-up year $t_0+\tau$, $\alpha_\tau$ are time fixed effects, $D_i$ is a dummy indicating divorce at $t_0$ and $\mathbf{x}_i$ are other variables measured prior to $t_0$. Just as in Equation \ref{eq:rq1_equation}, we estimate two specifications for Equation \ref{eq:rq2_equation_basic} by varying the inclusion of variables in  $\mathbf{x}_i$ based on absolute standardized mean differences after amtching. We complement the pooled hazard odds ratio with cumulative mortality risks using the Kaplan-Meier estimator (\citeauthor{kaplan1958nonparametric}, \citeyear{kaplan1958nonparametric}). We again report cumulative mortality risk among treated and controls, the risk difference per 1000 individuals and the corresponding risk ratio at the 5-year horizon\footnote{The hazard model uses up to ten years of follow-up although later follow-up years are available only for earlier divorce cohorts.}. We also run the lag-3 specification with more pre-divorce years.

In addition to the straightforward analysis of divorce on mortality, we run a mediation analysis to determine how much of the mortality difference can be attributed to `early' musculoskeletal rehabilitation as result of divorce. A single binary mediator measured one or two years after divorce is insufficient for this question since it can ignore deterioration that occurs later in the follow-up window, such as economic and rehabilitation histories that change after divorce. For example, divorce can potentially change earnings which can subsequently have the ability to affect the probability of entering musculoskeletal rehabilitation and dying. 

Let $D_i \in \{0,1\}$ is divorce status fixed at $t_0$, $\mathbf{z}_i$ the available pre-treatment matching variables for the relevant lag-1 and lag-3 specifications, $H=5$ the fixed mortality and mediation horizon, $\mathbf{L}_{i,t}$ the post-divorce history (\texttt{income}, \texttt{pension\_income}, \texttt{unsuccessful\_rehab\_application\_this\_year}, \texttt{non\_ms\_rehab}\footnote{A dummy indicating that an individual has a non-musculoskeletal rehabilitation that year.}, \texttt{income\_source}) in time $t$ and $M_{i,t-1} \in {0,1}$ is a dummy indicating whether individual $i$ entered musculoskeletal rehabilitation by the end of time $t-1$. We use information in time $t-1$ to predict information in time $t$. We cannot use information from time $t$ since we are unsure whether the dependent or independent variables are recorded first. We use weighted maximum likelihood to account for individuals actually observed in the data in a given year. 

We estimate the probability of an individual $i$ having a record for the year $t_o + t$ given that they are alive:

\begin{align}
\begin{split}
\text{Pr}(R_{i,t}=1 | D_i,\mathbf{z}_i, \mathbf{L}_{i,t-1}) = \Lambda_{\text{obs}}(\alpha_t + \beta D_i + \mathbf{z}_i'\psi + \mathbf{L}_{i,t-1}'\psi_L + \mu M_{i,t-1})
\end{split}
\tag{3}\label{eq:rq2_0_present}
\end{align}

We then predict $\hat{p}_{obs,i}(u)$, i.e. the probability that individual $i$ has a record in time $u$ given their information up until $u-1$ and inverse this to compute weight:

\begin{align}
\begin{split}
\hat{\omega}_{i,t} = \text{min} \left(\prod_{u=1}^t \frac{1}{\text{max}\left(\hat{p}_{\text{obs},i}(u),\underline{p}\right)}, \bar{\omega}\right)
\end{split}
\tag{4}\label{eq:rq2_0_weights}
\end{align}

We set the maximum weight, $\bar{\omega} = 20$ , to prevent one trajectory from having too much influence. We then estimate the parameters of a model on the probability of first musculoskeletal rehabilitation in year $t$ given that the individual did not experience musculoskeletal rehab in $t-1$ and other information in $t-1$\footnote{Weighed version: $\ell_{\text{rehab}} = \sum_{i,t} \hat{\omega}_{i,t} \left[M_{i,t}\text{log } \lambda_{i,t}^{\text{rehab}} + \left(1-M_{i,t} \right)\text{log } \left(1-\lambda_{i,t}^{\text{rehab}}\right)\right]$}:

\begin{align}
\begin{split}
\lambda_{i,t}^{\text{rehab}} &=
\text{Pr}(M_{i,t}=1 | M_{i,t-1}=0, D_i, \mathbf{z}_i, \mathbf{L}_{i,t-1}) \\ & = \Lambda_{\text{rehab}}(\kappa + \delta D_i + \mathbf{z_i}'\eta + \mathbf{L_{i,t-1}'}\eta_L )
\end{split}
\tag{5}\label{eq:rq2_1_unweighted}
\end{align}


Similarly, we estimate the parameters for a model on the probability of death in year $t$ given that an individual was alive at the end of $t-1$\footnote{Weighted version: $\ell_{\text{death}} = \sum_{i,t} \hat{\omega}_{i,t} \left[Y_{i,t}\text{log } \lambda_{i,t}^{\text{death}} + \left(1-Y_{i,t} \right)\text{log } \left(1-\lambda_{i,t}^{\text{death}}\right)\right]$} is given by:

\begin{align}
\begin{split}
\lambda_{i,t}^{\text{death}} & = \text{Pr}(Y_{i,t}=1| \text{alive at the end of } t-1) \\
& = \Lambda_{\text{death}}\left(\alpha_t + \theta_D D_i + \theta_M M_{i,t} + \theta_{DM} D_iM_{i,t} + \mathbf{z}_i' \rho + \mathbf{L}_{i,t-1}'\rho_L\right)
\end{split}
\tag{6}\label{eq:rq2_unweighted}
\end{align}

Since $\mathbf{L}_{i,t}$ is itself affected by divorce and evolves annually, we cannot hold it fixed at its observed values when generating the hypothetical divorce paths needed for the mediation decomposition\footnote{Suppose an individual earns 30,000 in 2013, divorces in 2014 and then earns 22,000 in 2015. Divorce could have caused the earnings to drop to 22,000 if the individual switched jobs. To determine the individual's rehabilitation or death risk if they had not divorced, we cannot plug in 22,000 because 22,000 is already corrupted by divorce. We need a clean number uncorrupted by divorce.}. Hence, we need to model the evolution of $\mathbf{L}_{i,t}$ among individuals who survive, i.e. observations where $Y_{i,t}=0$, so that we can simulate the counterfactual histories under the hypothetical divorce status for our mediation analysis (\citeauthor{robins1986}, \citeyear{robins1986}; \citeauthor*{lin2017survival}, \citeyear{lin2017survival}). We fit five separate models on the survivors, weighted by the same $\hat{\omega}_{i,t}$ (\citeauthor*{lin2017epi}, \citeyear{lin2017epi}) used in the weighted versions of Equations \ref{eq:rq2_1_unweighted} and \ref{eq:rq2_unweighted}. We use a log-scale linear regressions for \texttt{income} and \texttt{pension\_income}, logit for \texttt{unsuccessful\_rehab\_application\_this\_year} and \texttt{non\_ms\_rehab}, and multinomial logit for \texttt{income\_source}. These five models condition on the covariates used for matching $\mathbf{z}_i$, $\mathbf{L}_{i,t-1}$, current mediator $M_{i,t}$ and prior mediator $M_{i,t-1}$. The log-scale models for \texttt{income} and \texttt{pension\_income} additionally include \texttt{income\_source}, and the logit models include \texttt{income\_source}, \texttt{income} and \texttt{pension\_income}. Please see the section on \nameref{app:updating_history} in the Appendix. 

Once the parameters for these five equations are estimated, we simulate the counterfactual histories, one year at a time using the following process. 

\begin{enumerate}
  \item Rehab probability (uses $\mathbf{L}_{i,t-1}$)$\rightarrow$ draw$\rightarrow$ $M_{i,t}$
  \begin{itemize}
    \item Since our data starts from 2011, we use $\mathbf{L}_{i,2011}$ which contains true values to estimate the probability of rehabilitation, $\lambda_{i,2011}^{rehab}$. Since we need to feed $M_{i,2012}$ into the next step and not $\lambda_{i,2011}^{rehab}$, we draw a random number from a uniform distribution $U(0,1)$ and if it is smaller than $\lambda_{i,2011}^{rehab}$, then the value of $M_{i,2012}$ is set at 1 and 0 otherwise.
  \end{itemize}
  \item Death probability (uses $\mathbf{L}_{i,t-1}$, $M_{i,t}$) $\rightarrow$ if individual $i$ survives, continue
  \item History-update (uses $\mathbf{z}_i$, $D_i$, $\mathbf{L}_{i,t-1}$, $M_{i,t-1}$, $M_{i,t}$) $\rightarrow$ produces $\mathbf{L}_{i,t}$
  \begin{itemize}
    \item We use $\mathbf{L}_{i,2011}$, $M_{i,2011}$ and $M_{i,2012}$ to simulate $\mathbf{L}_{i,2012}$.
  \end{itemize}
  \item Hand $\mathbf{L}_{i,t}$ forward, repeat for $t+1$
\end{enumerate}


True divorce status $D_i$ is fixed at baseline and can simply be plugged into Equation \ref{eq:rq2_1_unweighted} for rehabilitation. In order to find the counterfactual effect, we can simply plug in the opposite value since there are only two divorce statuses and each status has a direct opposite. Hence, in step 1, we obtained the rehabilition probability and thus, predicted rehabilitation occurrence at time $t$, $M_{i,t}(d_i)$ for an individual with divorce status $d_i$. For individual $i$ whose plugged in divorce status is $d_i = 1$, denote the rehabilitation occurrence in time $t$ as $M_{i,t}(1)$ and when $d_i = 0$, $M_{i,t}(0)$. An individual's rehabilitation timing is given by $M_i(1)$ and $M_i(0)$\footnote{For example, if $M_i(d_i) = t+2$, then $M_{i,t'}(d_i)=0$ when $0<t'<t+2$ and $M_{i,t+2}(d_i)=1$.}.We then obtain the sets, $M(1)$ and $M(0)$, which contain the rehabilitation timing for all individuals. 

Subsequently, when we cross $d_i \in \{0,1\}$ with the resulting $M_{i}(d_i)$ and $M_i(d_i')$ where $d_i \neq d_i'$, we yield four different possibilities $\in \left\{(0,M_{i}(0)), (0,M_{i}(1)), (1, M_{i}(0)), (1, M_{i}(1))\right\}$. We might be tempted to plug in these four possibilities into Equations \ref{eq:rq2_1_unweighted}, \ref{eq:rq2_unweighted} and the equations for updating $\mathbf{L}_{i,t}$ (in the Appendix) but $M_i(d_i)$ is not a fixed characteristic but the output of a year-by-year stochastic process. In other words, since $M_{i}(d_i)$ is based on \textit{at least one random} draw under a hypothetical divorce status, there is no direct opposite like in the case with divorce status. Furthermore, we have two dimensions, $d_i$ and $M_{i}(d_i)$, where we need a counterfactual but it is difficult to defend \textit{individual} counterfactuals for $(0,M_{i}(1))$ and $(1, M_{i}(0))$. Think of this like a natural experiment. If we only wanted a clean effect of divorce, we can externally impose divorce on some individuals. The individuals who we forced to divorce would experience the rehabilitation timing of someone with a divorced status, likewise for the non-divorced. If we wanted a clean direct effect of divorce and an indirect effect attributable to rehabilitation, we not only need to force someone to divorce but also experience the rehabilitation timing of a non-divorced person at the same time and vice versa, i.e. we need the off-diagonal or counterfactuals. This is not possible with a natural experiment as we cannot impose both dimensions simultaneously. For example, suppose we want to plug in the $(0,M_{i}(1))$ combination. Suppose the probability of rehabilitation under divorce is 80\% and under non-divorce 5\%. This means the timing of $M_i(1)$ is likely to be earlier than that of $M_i(0)$. Hence, plugging in $d_i = 0$ together with $M_i(1)$ (which is, to put it crudely, `contaminated' with individual $i$'s identity when $d_i = 1$) into Equation \ref{eq:rq2_unweighted} would be like imposing both dimensions simultaneously. This argument also holds for the $(1, M_{i}(0))$ combination. To manage this `contamination', we sever the link between individual $i$ and their estimated rehabilitation timing by attaching another individual's rehabilitation timing. While $(0, M_i(0))$ and $(1, M_i(1))$ do not have this `contamination' problem, we also sever the link between the individual and estimated rehabilitation timing for consistency.

To do this, we make use of the distribution of $M(1)$ and $M(0)$. This would then be similar to imposing an external rehabilitation timing. Before step 2, we replace individual $i$'s value of $M_i(d_i)$ with individual $j$'s $M_j(d_j)$ where $d_i = d_j$. This means that when both individuals $i$ and $j$ have an identical plugged in divorce status, individual $i$ will have individual $j$'s rehabilitation timing. This second-best approach helps us identify what the average death risk for individuals with a particular plugged in divorce status looks like under a random redistribution of rehabilitation timing\footnote{One caveat is that, it may be possible that individual $j$'s rehabilitation timing is unlikely to be produced given individual $i$'s rehabilitation probability, for example, if individual $j$'s probability of rehabilitation in year 1 is high (e.g., 90\%) while individual $i$'s is low (e.g., 10\%). This is called a positivity or overlap concern (\citeauthor*{vanderweele2014effect}, \citeyear{vanderweele2014effect}; \citeauthor*{petersen2012diagnosing}, \citeyear{petersen2012diagnosing}).}.

Suppose individual $i$ received individual $j$'s value when $d_i = 0$, and individual $k$'s value when $d_i = 1$. This means that individual $i$'s hypothetical rehabilitation timings going into step 2 take the values of $M_j(0)$ and $M_k(1)$ respectively. We can then plug in the four combinations, $\in \left\{(0,M_{j}(0)), (0,M_{k}(1)), (1, M_{j}(0)), (1, M_{k}(1))\right\}$, into Equation \ref{eq:rq2_unweighted} in step 2 for individual $i$ in time $t$. We then update $\mathbf{L}_{i,t}$  and repeat all four steps until we reach five years. 

Since individual $i$'s $M_{i,t}$ depends on a random draw which subsequently affects $L_{i,t}$, we re-run the entire process 25 times with different draws for the rehabilitation occurrence to average out simulation randomness. Let $M_{i,r}(d_i)$ and $Y_{i,r}(d_i,M(d_i'))$ denote individual $i$'s simulated rehabilitation and death indicators by the horizon of five years from $t_0$ and $R=25$ denote the number of Monte Carlo repeats. The Monte Carlo average for individual $i$ is given by:

\begin{align}
\begin{split}
\hat{M}_{i}(d_i) & = \frac{1}{R} \sum_{r=1}^{R}M_{i,r}(d_i) \\
\hat{Y}_i(d_i,M(d_i')) & = \frac{1}{R} \sum_{r=1}^R Y_{i,r}(d_i,M(d_i'))
\end{split}
\tag{7}\label{eq:avg_within_individual}
\end{align}

We then average across the matched treated individuals. Let $\mathcal{T}$ denote the set of $N_M$ matched treated individuals:

\begin{align}
\begin{split}
\hat{M}(d) & = \frac{1}{N_M} \sum_{i \in \mathcal{T}} \hat{M}_i(d_i) \\
\hat{Y}(d, M(d')) & = \frac{1}{N_M} \sum_{i \in \mathcal{T}} \hat{Y}_i(d_i,M(d_i'))
\end{split}
\tag{8}\label{eq:avg_between_individual}
\end{align}

In Equation \ref{eq:avg_between_individual}, we plug in the four different combinations to obtain \textit{population} averages of the rehabilitation risk under the hypothetical divorce status, $d$ and mortality risk under the hypothetical divorce $d$ combined with the rehabilitation timing of divorce status $d'$ where $d, d' \in \{0,1\}$. We can then decompose the total effect into the direct effect and indirect effect:

\begin{align}
\begin{split}
\widehat{\text{Total effect}}&= \hat{Y}(1,M(1))-\hat{Y}(0,M(0)) \\
\widehat{\text{Direct effect}}&= \hat{Y}(1,M(0))-\hat{Y}(0,M(0)) \\
\widehat{\text{Indirect effect}}&= \hat{Y}(1,M(1))-\hat{Y}(1,M(0))
\end{split}
\tag{9}\label{eq:decomposition}
\end{align}

The total effect is simply answering what the difference in mortality risk is between divorced and non-divorced individuals within five years. The direct effect captures part of the difference in mortality risk related to divorce but unrelated to rehabilitation, i.e. given that rehabilitation timing is held constant, how does divorce (potentially operating through financial strain, loneliness, stress) affect mortality risk. The indirect effect answers how much of the difference in mortality risk is attributed to divorce shifting individuals onto the rehabilitation timing distribution typical of divorced individuals. We calculated interventional direct and indirect effects (\citeauthor*{vanderweele2014effect}, \citeyear{vanderweele2014effect}) which is different from natural direct and indirect effects  (\citeauthor{pearl2001direct}, \citeyear{pearl2001direct}). A natural indirect effect answers what would happen to individual $i$'s mortality risk if their rehabilitation timing shifted from what it looks like when individual $i$ is non-divorced to what it would look like if individual $i$ is divorced while holding divorce status constant at $d=1$. In contrast, an interventional indirect effect answers what would happen to average mortality risk if the population's rehabilitation timing shifts from the distribution typical of non-divorced individuals to a distribution typical of divorced individuals while holding divorce status constant at $d=1$. 

Finally, we then bootstrap the sample 200 times (\citeauthor{efron1993bootstrap}, \citeyear{efron1993bootstrap}) with replacement and repeat the entire procedure. We do this to obtain the 95\% confidence interval for the total, direct and indirect effects. 




\section{Results}
\subsection{Descriptive summary}

{
\color{red}

\begin{itemize}
\item How many individuals are in the 20\% dataset? How many person years?
\item How many individuals have undergone divorce each year?
\item Matching limited to 2012 to 2018 for lag1 spec, and 2014 to 2018 for lag3 spec. Who do we filter out? Sample size and number of matches used for each. Do people have gaps in the data?
\item How many people undergo medical rehabilitation? What are the top rehabilitation codes and cases? What happens in the five year follow up? Do people die before they get the rehab? How many complete the rehab and survive the 5 years? How many neither get rehab nor die in the 5 years?
\item $N_T$: eligible treated individuals, $N_M$: matched individuals (treated and controls), gender composition. Are there people who show up twice, perhaps first as control and then as divorced down the line?
\item What is the matching rate by treatment year?
\item Distribution of propensity scores, mean, max between matched individuals
\item Event study coefficients, pre-trend Wald test/joint coefficient test
\end{itemize}
}


\subsection{RQ1: Musculoskeletal rehabilitation}

{
\color{red}
\begin{itemize}
\item How much more likely are divorced people to start rehab than almost identical non-divorced people?
\item Does adding leftover covariates which were not well balanced affect the above result?
\item Does the risk mostly show up in the year following divorce or does it creep up by year 5?
\item What is the share of divorced people who started rehab by year 5 compared to controls?
\item What is the difference in the results when lag1 spec is used versus lag3 spec?
\item What is the difference in rehab risk between genders?
\item Of the matches, how many have a full five years of usable follow up data?
\end{itemize}
}

\subsection{RQ2: Mortality}

{
\color{red}

\begin{itemize}
\item How much more likely are divorced people to die during the follow up period
\item Compare lag1 and lag3 spec mortality risk
\item Total, direct and indirect effects plus 95\% confidence interval.
\item How often is it the case that the borrowed rehab timing is completely off?
\item Total, direct and indirect effects for men vs women. 
\end{itemize}

}

\section{Discussion}

\section{Conclusion}

\printbibliography[title={References}]
% \bibliographystyle{apalike}
% \bibliography{main}

\section*{Appendix}
\subsection*{Rehabilitation system in Germany}
\phantomsection
\makeatletter
\def\@currentlabelname{Rehabilitation system in Germany}
\makeatother
\label{app:rehab_germany}
For an individual covered by statutory health insurance, musculoskeletal pain first enters the health insurance system through sickness absence rather than through the DRV. If pain makes an employee unable to work, the individual must immediately inform the employer of the incapacity and the expected duration. If the incapacity lasts longer than three calendar days, a medical certificate needs to be acquired from the individual's general practitioner, although the employer may require certification from the first day. The employer continues to pay the normal salary for up to six weeks for incapacity caused by the same illness (\citeauthor{entgfg1994}, \citeyear{entgfg1994}). After the six weeks, if the illness continues, the statutory health insurer steps in and pays \textit{Krankengeld}, which is usually about 70\% of covered gross income up to a maximum of 90\% of net income. For the same illness, this benefit is limited to 78 weeks within a 3-year period, counted from the first day of incapacity, meaning that the time of illness paid by the employer is included in the 78 weeks (\citeauthor{sgbv1988}, \citeyear{sgbv1988}). Medical rehabilitation facilitated by the DRV requires a separate application and assessment of whether earning capacity is substantially endangered or already reduced, if rehabilitation is expected to prevent or improve that reduction (\citeauthor{sgbvi1989}, \citeyear{sgbvi1989}; \citeauthor{drvGeneralReha}, \citeyear{drvGeneralReha}). This application is usually independent of the six weeks and \textit{Krankengeld} although the health insurer can allow the individual ten weeks to submit a rehabilitation application indicating that earning capacity is threatened (\citeauthor{sgbv1988}, \citeyear{sgbv1988}). For musculoskeletal conditions, the DRV assessment focuses on documented functional limitations and their consequences for work and social participation. Chronic pain, psychosocial circumstances surrounding incapacity, comorbidity, and active cooperation are considerations in this assessment rather than four cumulative eligibility conditions. Isolated pain is insufficient to establish rehabilitation need for certain hip, knee, and shoulder conditions (\citeauthor{drvMskCriteria2005}, \citeyear{drvMskCriteria2005}). Approved medical rehabilitation is normally provided on an inpatient or full-day outpatient basis and usually lasts approximately three weeks. The employer continues wages during rehabilitation if entitlement remains, otherwise, the DRV may pay \textit{Ubergangsgeld}, normally 68\% of previous net earnings or 75\% for an insured person with a qualifying child, while entitlement to \textit{Krankengeld} is suspended (\citeauthor{drvUebergangsgeld,sgbv1988}, \citeyear{drvUebergangsgeld,sgbv1988}).


\subsection*{Dealing with a large sample of controls}
\phantomsection
\makeatletter
\def\@currentlabelname{Dealing with a large sample of controls}
\makeatother
\label{app:control_limitation_500_000}

Suppose in 2015, there are 50,000 individuals who are divorced and 2 million eligible controls. Using all 2 million controls to estimate the logit model would cause memory issues and take a long time to run. So we imposed a cap of $450,000 (500,000 - 50,000)$ eligible controls which are then randomly selected, although the same seed is used for reproducibility. If there are less than 450,000 eligible controls, we use all of them to estimate the parameters for the logit model. Once we have estimated the logit model, we estimate the propensity score for everyone, not only the 450,000 controls and 50,000 treated. Then we match based on propensity scores and end up with a treated-control pair, i.e. in 2015, we would have 50,000 treated and 50,000 controls.

\subsection*{Updating the evolution of $\mathbf{L}_{i,t}$}
\phantomsection
\makeatletter
\def\@currentlabelname{Updating the evolution of $\mathbf{L}_{i,t}$}
\makeatother
\label{app:updating_history}

We first estimate the parameters of the model for \texttt{income\_source}. Let \texttt{income\_source} be denoted by $S_{i,t}$:

\begin{align*}
\text{Pr}(S_{i,t} = k | \cdot) = \frac{\exp \left(c_k+\delta_kD_i +\mathbf{z}_i'\eta_k + \mathbf{L}_{i,t-1}'\eta_{L,k} + \mu_kM_{i,t-1} + \zeta_k M_{i,t} + \tau_{k,t}\right)}{\sum_{k'} \exp(\cdot)}
\end{align*}

where $c_k$ is the constant and $\tau_{k,t}$ is the time fixed effect at year $t$ for category $k$. 

Then, we estimate the parameters for the models on \texttt{income} and \texttt{pension\_income} using the equation structure below. Let $A_{i,t}$ represent \texttt{income} and \texttt{pension\_income}:

\begin{align*}
\log(1+A_{i,t}^v) = c_v + \beta_v D_i + \mathbf{z}_i'\gamma_v + \mathbf{L}_{i,t-1}'\gamma_{L,v} + \mu_v M_{i,t-1} + \zeta_v M_{i,t} + \tau_{v,t} + \phi_v' \mathbbm{1}[\tilde S_{i,t}] + \varepsilon_{i,t}^v
\end{align*}

where $v$ indicates whether we are estimating \texttt{income} or \texttt{pension\_income}, $c_v$ is the constant, $\tilde S_{i,t}$ was the previously estimated \texttt{income\_source} and $\tau_{v,t}$ is the time fixed effect. 

We finally estimate the parameters for the model on the binary indicators \texttt{non\_ms\_rehab} and \texttt{unsuccessful\_rehab\_application\_this\_year}:

\begin{align*}
\Pr\!\left(B_{i,t}^{1}=1 \mid \cdot\right)
&=
\Lambda\!\Big(
    \alpha_1
    + \theta_1 D_i
    + \mathbf{z}_i'\rho_1
    + \mathbf{L}_{i,t-1}'\rho_{L,1}
    + \mu_1 M_{i,t-1}
    + \nu_1 M_{i,t} \\
&\qquad
    + \tau_{1,t}
    + \widetilde{\mathbf{A}}_{i,t}'\pi_1
    + \mathbbm{1}\!\left[\widetilde{S}_{i,t}=1\right]\pi_{S,1}
\Big)
\end{align*}

\begin{align*}
\Pr\!\left(B_{i,t}^{2}=1 \mid \cdot,\ B_{i,t}^{1}\right)
&=
\Lambda\!\Big(
    \alpha_2
    + \theta_2 D_i
    + \mathbf{z}_i'\rho_2
    + \mathbf{L}_{i,t-1}'\rho_{L,2}
    + \mu_2 M_{i,t-1}
    + \nu_2 M_{i,t} \\
&\qquad
    + \tau_{2,t}
    + \widetilde{\mathbf{A}}_{i,t}'\pi_2
    + \mathbbm{1}\!\left[\widetilde{S}_{i,t}\right]\pi_{S,2}
    + \pi_B \widetilde{B}_{i,t}^{1}
\Big)
\end{align*}


where $B_{i,t}^1$ is a dummy representing \texttt{unsuccessful\_rehab\_application\_this\_year} and is estimated first, $B_{i,t}^2$ is a a dummy representing \texttt{non\_ms\_rehab} which is estimated second hence also contains the predicted $\widetilde{B}_{i,t}^{1}$ and $\widetilde{\mathbf{A}}_{i,t}$ is a vector containing the drawn level (not $\log$) amounts of (\texttt{income } \texttt{pension\_income}) which we just estimated in the log-scale linear regressions above.

\end{document}
